% Owned proposed fragment B.123; no central register mutation.
\subsection*{Two actual source operators and the complete unfilled kernel}

\statement{B.123}{Proved}{The whole ordinary one-neck group has an exact free-eleven, two-source-operator presentation, and its fixed-source question reduces to the actual family kernel in both categories}
\label{stmt:B123}
Use the actual common marking of B.114 and B.119:
\[
 K=F(a,B,d,e,f),\qquad c=c_-,\qquad \tau=c_+,\qquad
 m=k_-,\quad p=k_+.
\]
The two cones are odd and the other generators are ordinary.
Let
\begin{equation}\label{eq:B123-generic-fiber}
 E_0=K*C_2(c)*C_2(\tau),\qquad
 F=\ker(E_0\longrightarrow C_2).
\end{equation}
Then $F$ is the actual free group of rank eleven, with marked basis
\begin{equation}\label{eq:B123-free-eleven}
 a,B,d,e,f,\ cac,cBc,cdc,cec,g=cfc,\ x=\tau c.
\end{equation}
There are two explicitly specified automorphisms $\alpha_m,\alpha_v$
of $F$, defined below by the complete regional markings and partial
inner automorphisms. There is a finite list $\mathcal A$ of
\emph{forty} marked words, namely both ordinary sheets of all twenty
regional node words. Put
\begin{equation}\label{eq:B123-whole-closure}
 \begin{gathered}
 S=F(m,v),\qquad
 N=\left\langle\!\left\langle
       \alpha_s(r):r\in\mathcal A,\ s\in S
       \right\rangle\!\right\rangle_F,\qquad W=F/N.
 \end{gathered}
\end{equation}
Here $s$ ranges over \emph{every signed word} in $m,v$, with
$\alpha_s(w)=sws^{-1}$; no bound on its length is imposed.
There is an exact marked isomorphism
\begin{equation}\label{eq:B123-actual-splitting}
 \mathcal G^0_{\rm ann}=W\rtimes S.
\end{equation}
The actual fifth source word is $v$, and its sixth source word is
$m$. The four contact source words give explicit elements
$s_1,\ldots,s_4\in W$. Their full image is
\begin{equation}\label{eq:B123-source-orbit}
 H=I\rtimes S,\qquad
 I=\left\langle\alpha_s(s_i):s\in S,\ 1\leq i\leq4\right\rangle
       \leq W.
\end{equation}
The brackets defining $I$ denote a subgroup, not an ambient normal
closure.

Conjugation $\delta=\operatorname{Inn}(\tau)$ preserves $W$ and
fixes $S$ pointwise. Let $J=\ker(W\to F(f,g,x))$ be the actual
eight-family filling kernel. Then
\begin{equation}\label{eq:B123-discrete-gate}
 (\mathcal G^0_{\rm ann})^\delta=W^\delta\rtimes S,
 \qquad W^\delta=J^\delta,\qquad
 (\mathcal G^0_{\rm ann})^\delta=H
       \ \Longleftrightarrow\ J^\delta=I.
\end{equation}
There are also exact full-completion comparisons
\begin{equation}\label{eq:B123-compact-gate}
 \begin{gathered}
 \widehat{\mathcal G^0_{\rm ann}}=\widehat W\rtimes\widehat S,
 \qquad \widehat W=\widehat F/\overline N,\qquad
 \overline H=\overline I\rtimes\widehat S,\\
 \bigl(\widehat{\mathcal G^0_{\rm ann}}\bigr)^\delta
       =\widehat W^{\,\delta}\rtimes\widehat S,
 \qquad \widehat W^{\,\delta}=\overline J^{\,\delta},\\
 \bigl(\widehat{\mathcal G^0_{\rm ann}}\bigr)^\delta=\overline H
       \ \Longleftrightarrow\ \overline J^{\,\delta}=\overline I,
 \qquad
 \overline J=\ker\bigl(\widehat W\to\widehat{F(f,g,x)}\bigr).
 \end{gathered}
\end{equation}
Here $\overline N$ is the closure of $N$ in $\widehat F$;
$\overline I$ is the closure of the discrete actual source-orbit
subgroup in $\widehat W$; and $\overline H$ is the image of the
completed actual source. The notation $\overline J$ means the
displayed actual closed kernel. No identification of it with
$\widehat J$, or of $\overline N$ with $\widehat N$, is asserted.
The equalities $J^\delta=I$ and
$\overline J^{\,\delta}=\overline I$ remain open.

\begin{proof}
\emph{B.123.1. Both complete regional free-product bases.}
The actual physical neck basis and mirror are
\begin{equation}\label{eq:B123-eta}
 \begin{gathered}
 (O_1,O_2,O_3,O_4,O_5)
  =(aBa^{-1},a,fde d^{-1}f^{-1},d,d^{-1}fd),\\
 P_i=O_1\cdots O_i,\quad P_0=1,\qquad
 \eta(O_i)=P_{i-1}O_i^{-1}P_{i-1}^{-1},\quad \eta^2=1.
 \end{gathered}
The inverse basis change is
\[
 a=O_2,\quad B=O_2^{-1}O_1O_2,\quad d=O_4,\quad
 f=O_4O_5O_4^{-1},\quad e=O_5^{-1}O_4^{-1}O_3O_4O_5.
\]
These specify $\eta$ on the full common $K$ without a new
peripheral conjugation convention. Define
\begin{equation}\label{eq:B123-local-tuples}
 \begin{gathered}
 (a_-,B_-,d_-,e_-,f_-)=(a,B,d,e,f),\\
 (a_+,B_+,d_+,e_+,f_+)=(\eta(a),\eta(B),\eta(d),\eta(e),\eta(f)),\\
 C_-=a^{-1}ca,\qquad C_+=a_+^{-1}\tau a_+,\qquad
 b_\pm=C_\pm B_\pm C_\pm.
 \end{gathered}
The two ordered systems
\[
 (a_-,b_-,d_-,e_-,f_-,C_-,\tau),\qquad
 (a_+,b_+,d_+,e_+,f_+,C_+,c)
\]
are free-product bases of $E_0$. For the minus system the inverse
change sends $B$ to $C_-b_-C_-$ and $c$ to $a_-C_-a_-^{-1}$.
For the plus system first recover its ordinary $B_+$ as
$C_+b_+C_+$, apply $\eta^{-1}=\eta$ to its five ordinary
coordinates, and recover $\tau=a_+C_+a_+^{-1}$, keeping $c$.
Thus both inverses are explicit; the other cone is part of each
complete basis, not an omitted peripheral label.

\emph{B.123.2. Exact generic monodromy, in both directions.}
In either local system abbreviate the corresponding far lift by
$k$, its cone by $C$, and its branch by $f$. The six actual generic
regional relations are
\begin{equation}\label{eq:B123-generic-relations}
 [h,C]=1,\qquad hkh^{-1}=CkC,\qquad
 [k,U]=1\quad(U=a,b,d,e),\qquad h=fk.
\end{equation}
They are equivalent to
\begin{equation}\label{eq:B123-partial-inner}
 kCk^{-1}=f^{-1}Cf,\qquad
 kfk^{-1}=f^{-1}CfCf,\qquad kUk^{-1}=U.
\end{equation}
The first equation is exactly $[fk,C]=1$. The other generic
equation says $fkf^{-1}=CkC$; substituting the first equation
and solving for $kfk^{-1}$ gives the second displayed image.
Conversely those two images recover both equations in
\eqref{eq:B123-generic-relations}. These are precisely the
partial inner automorphism $\operatorname{Inn}(f^{-1}C)$ on
$\langle C,f\rangle$, fixing the four outsiders. Its inverse
is the partial $\operatorname{Inn}(Cf)$.

In the minus basis extend this automorphism by fixing $\tau$;
in the plus basis extend it by fixing $c$. These extensions are
exactly the two actual cross relations $[\tau,m]=[c,p]=1$.
Transport the two maps and inverses through
\eqref{eq:B123-local-tuples}; call the resulting automorphisms
of $E_0$ $\alpha_m,\alpha_p$. They preserve ordinary parity.
Since there is no relation between the two base letters, the
complete generic-plus-cross presentation is exactly
\begin{equation}\label{eq:B123-generic-product}
 \mathcal G_{\rm gen}=E_0\rtimes F(m,p).
\end{equation}
The explicit converse just proved gives inverse presentation maps,
so this is more than a map onto a semidirect product. It is the
actual two-loop regular-spine presentation before the node disks
are attached; no bundle over the critical total base is assumed.

\emph{B.123.3. All twenty node words and both ordinary sheets.}
After \eqref{eq:B123-generic-relations}, the ten actual node
words in a local tuple $(a,b,d,e,f,C)$ are exactly
\begin{equation}\label{eq:B123-ten-nodes}
 \begin{gathered}
 [U,CUC]\quad(U=a,b,d,e),\\
 [a,b],\quad[CaC,d],\quad[CaC,e],\quad[d,CeC],\\
 [b,fdf^{-1}],\quad[b,fef^{-1}].
 \end{gathered}
For the last two, the original words use $hdh^{-1}$ and
$heh^{-1}$; $h=fk$ and the actual outsider equations replace
them by $fdf^{-1}$ and $fef^{-1}$. Substitute both complete
tuples \eqref{eq:B123-local-tuples} into
\eqref{eq:B123-ten-nodes}. This specifies all twenty marked
orbifold words $r$, not just their conjugacy classes or incidence.

Use the ordinary Schreier transversal $(1,c)$ in $E_0$.
An even letter $u^{\pm1}$ from $K$ contributes $u^{\pm1}$
in state zero and $(cuc)^{\pm1}$ in state one; $c$ changes
state without a letter; $\tau$ changes state and contributes
$x$ in state zero, $x^{-1}$ in state one. This gives precisely
the free basis \eqref{eq:B123-free-eleven}, with inverse
expansions $cuc$ and $x=\tau c$. For every even node word take
the two rewrites representing $r$ and $crc$. These forty words
are the list $\mathcal A$ in \eqref{eq:B123-whole-closure}.
If they are denoted $(r_0,r_1)$, then literally
\begin{equation}\label{eq:B123-node-cone-mates}
 \delta(r_0)=x r_1x^{-1},\qquad
 \delta(r_1)=x r_0x^{-1}.
\end{equation}
The cone action on the whole free basis is
\begin{equation}\label{eq:B123-delta}
 \delta(u)=x(cuc)x^{-1},\qquad
 \delta(cuc)=xux^{-1}\quad(u\in\{a,B,d,e,f\}),\qquad
 \delta(x)=x^{-1}.
\end{equation}
It squares to the identity freely.

\emph{B.123.4. The actual fixed source base.}
Put $A=a_+$ and
\begin{equation}\label{eq:B123-source-base}
 q=Af_+A^{-1},\qquad
 v=Af_+pA^{-1}=qp.
\end{equation}
The last equality uses $\alpha_p(A)=A$, one of the exact
outsider equations. The displayed $v$ is the unchanged fifth
raw source word of B.119, including its entire gap conjugation;
the sixth is $m$. It centralizes $\tau$ already in
\eqref{eq:B123-generic-product}, because it is the whole
$A$-conjugate of $h_+=f_+p$, which centralizes $C_+$.
The other fixed equation is the actual $[\tau,m]=1$.
Consequently the induced source automorphisms are
\begin{equation}\label{eq:B123-source-actions}
 \alpha_m|_F,\qquad
 \alpha_v=\operatorname{Inn}(q)\alpha_p|_F,\qquad
 \alpha_v^{-1}=\alpha_p^{-1}\operatorname{Inn}(q^{-1})|_F,
\end{equation}
and all their signed versions commute with $\delta$.

The first four raw source words are also fixed unambiguously.
In the plus tuple put
\[
 E_+=f_+d_+e_+d_+^{-1}f_+^{-1},\qquad
 T_+=C_+b_+C_+b_+C_+E_+C_+b_+^{-1}C_+b_+^{-1},
\]
and $\Delta_U=UC_+UC_+$. Their words in $F$ are
\begin{equation}\label{eq:B123-contact-words}
 \begin{gathered}
 s_1=A(C_+b_+C_+b_+)A^{-1},\qquad s_2=AT_+^2A^{-1},\\
 s_3=A\bigl(C_+b_+[d_+,C_+d_+C_+]b_+^{-1}C_+\Delta_{d_+}\bigr)A^{-1},\\
 s_4=A\bigl(C_+b_+[a_+,C_+a_+C_+]b_+^{-1}C_+\Delta_{a_+}\bigr)A^{-1}.
 \end{gathered}
\end{equation}
These are the literal raw source forms of B.112 and B.119,
rewritten by the same ordinary rule. Their images become
$\delta$-fixed after the actual node relations; their free
words need not be fixed before those disks are attached.

The map $F(m,v)\to\mathcal G_{\rm gen}$ is injective: its
projection to the free base sends $m,v$ to the free basis $m,p$.
It is therefore an actual source splitting. Replacing $p$ by
$q^{-1}v$ changes its fiber automorphism only by the inner
automorphism $\operatorname{Inn}(q^{-1})$.

\emph{B.123.5. The whole normal closure and the exact source image.}
The ordinary normal closure of all orbifold node disks is generated
by their two cosets $1,c$. In the ordinary generic group, further
conjugations give precisely the $F$-conjugates and all signed
base translates of the forty words. For a subgroup normal in $F$,
invariance under $\alpha_m,\alpha_p$ and their inverses is
equivalent to invariance under $\alpha_m,\alpha_v$ and their
inverses, by \eqref{eq:B123-source-actions}. Thus the actual
normal subgroup killed is exactly $N$ in
\eqref{eq:B123-whole-closure}, inside the fiber. This proves
\eqref{eq:B123-actual-splitting}. In particular the node filling
does not impose a relation on $S$. Equations
\eqref{eq:B123-node-cone-mates} and the source commutation
prove that $N$ is $\delta$-stable.

For an additional complete marked comparison, send the ordinary
thirteen-generator system of B.119 into this normal form by keeping
the ten ordinary puncture generators, sending its $x$ to the
eleventh fiber generator, $m$ to $m$, and $p$ to $q^{-1}v$.
All twenty-four generic ordinary rows become identities, and
all other forty rows become exactly the forty specified node
words. Both sheets and their original names are retained. This
is an executed check of the entire generating scheme; the
normal-closure argument supplies all its signed-orbit consequences.

The actual source has the six based free generators
$(s_1,s_2,s_3,s_4,v,m)$ before taking its image. Filling its
four contact punctures is the literal free retraction
$F_6=F_4*F_2\to F_2$. Its kernel is generated as a subgroup
by all conjugates of the four contact generators under the
free base. Taking the actual image yields exactly
\eqref{eq:B123-source-orbit}, with $H\cap W=I$.
The source geometry, or the raw local source identities of
B.112, gives $I\leq W^\delta$. Since $\delta$ fixes $S$
pointwise, unique semidirect decomposition proves the first
and last comparisons in \eqref{eq:B123-discrete-gate}, with
$W^\delta$ initially in place of $J^\delta$.

\emph{B.123.6. The actual family filling.}
On \eqref{eq:B123-free-eleven} kill $a,B,d,e$ and their
four $c$-mates, and keep $f,g,x$. Every word in
$\mathcal A$ becomes the identity. The four signed source
actions and $\delta$ descend literally to the corresponding
actions in B.119; in particular $q\mapsto f^{-1}$ and
$v\mapsto f^{-1}p$. This gives the actual map $W\to F(f,g,x)$.
The exact ordinary-fiber fixed subgroup in that quotient is
trivial, by B.119. Therefore every element of $W^\delta$ lies
in its kernel $J$. The reverse inclusion of fixed subgroups is
tautological, proving $W^\delta=J^\delta$ and completing
\eqref{eq:B123-discrete-gate}.

\emph{B.123.7. Full topology and actual closed kernels.}
The group $W$ is finitely generated. Its characteristic finite-index
subgroups are cofinal: intersect the finitely many subgroups of
index at most any prescribed finite index. Given a finite target
of $W\rtimes S$, choose such a subgroup $U\lhd W$ inside its
restriction kernel, and choose $V\lhd S$ of finite index inside
both its restriction kernel and the action kernel on $W/U$.
Then
\[
 (W\rtimes S)\longrightarrow(W/U)\rtimes(S/V)
\]
has kernel contained in that prescribed target kernel. These
finite targets are cofinal. They also refine every finite target
on $W$, while $S$ is a retract. Hence both factors have their
full induced topology and completion gives the exact split
product in \eqref{eq:B123-compact-gate}. This uses finite
generation, not residual finiteness of $W$.

The universal property on finite targets gives
$\widehat W=\widehat F/\overline N$. Each characteristic
$U$ is $\delta$-stable, and $\delta$ fixes $S$ pointwise.
Thus unique decomposition in the completed split product gives
its fixed subgroup exactly as stated.

The image of the completed actual source is the closure of $H$.
Closing $I\rtimes S$ in this split product gives precisely
$\overline I\rtimes\widehat S$: the closed subgroup
$\overline I$ is preserved by $\widehat S$, by continuity,
and projection to the two factors shows both inclusions.
Equivalently the completed free-retraction kernel of the source
is the closure of its discrete kernel, whose actual image is
$\overline I$. No abstract completion of $I$ is used.

Finally the family-filling diagrams are equivariant. B.119 gives
trivial fixed subgroup in the completed ordinary fiber
$\widehat{F(f,g,x)}$. Hence
$\widehat W^{\,\delta}=\overline J^{\,\delta}$ for the
actual closed kernel defined in the statement. This proves
all of \eqref{eq:B123-compact-gate}. The base source splitting
is actual; no section of the entire coarse group into the
family-unfilled group is inferred.
\end{proof}

\paragraph{Complete marked verification.}
The producer reads only the two stable actual presentation JSONs.
It checks both regional basis inverses, both signed far actions,
all generic relations, the actual source correction and all
forty-four source/involution comparisons, all forty node words,
all sixty-four original ordinary rows, and all six raw source
words. Its companion verifies the forty-four signed operator
filling diagrams, all eleven cone diagrams, and the vanishing
of all forty node and four contact words. Neither imports a
quotient oracle or asserts that a finite orbit closes. The exact
complete operator closure in \eqref{eq:B123-whole-closure}
remains part of the theorem.
\begingroup\raggedright
\path{work_CE/code/compile_annular_two_operator25.py},
\path{work_CE/code/compile_annular_two_operator25.out},
\path{work_CE/next_symmetric_nodes/annular_two_operator25.json};
\path{work_CE/code/check_two_operator_filling25.py},
\path{work_CE/code/check_two_operator_filling25.out},
\path{work_CE/next_symmetric_nodes/two_operator_filling25.json};
\path{work_CE/annular_two_operator_reduction25.txt}.
\par\endgroup
